\chapter{Model Configurations and Implementation Details}
\label{app:implementation}

\section{Executed Component Boundaries}
The PHM scalar-image regressor and the REASSEMBLE action classifier use different
components of RT-DETR. Table~\ref{tab:component-boundaries} identifies the model
weights, intermediate representation and output used by each branch.

\begin{table}[htbp]
\caption{Inputs and Outputs of the Executed Model Components}\label{tab:component-boundaries}
\centering\small

\begin{tabular}{P{0.19\linewidth}P{0.34\linewidth}P{0.37\linewidth}}\toprule
Component & Input and representation & Executed boundary\\\midrule
PHM scalar image branch & Three spatially pooled 256-dimensional maps, plus recorded
extraction latency: 769 actual head covariates & Ultralytics RT-DETR-L prediction path
with hooked maps; not the CNN-only REASSEMBLE path\\
REASSEMBLE visual branch & Sixteen RGB frames pooled to an 896-dimensional action
descriptor & Frozen ResNet18-vd CNN backbone inside RT-DETR-R18vd, followed by a trained
binary head\\
REASSEMBLE numerical statistics & 22 channels, ten statistics each: 220 features &
Training-only standardisation and logistic regression\\
REASSEMBLE temporal branch & 512 progress positions by 22 channels & Scratch-initialised
compact PatchTST and fitted projection/head\\
External numerical adapters & MSM: 90; HANDOVER: 36; ViFailback: 280 statistics &
Dataset-specific input construction; not an unchanged raw-sequence temporal
encoder\\\bottomrule
\end{tabular}
\displaynote{The reported dimensions apply to the listed implementation paths.}
\end{table}
\ev{Executed component map}{\EvidenceRef{E108}{}}

The REASSEMBLE extractor accesses the CNN through
\nolinkurl{wrapper.model.backbone.model}; the detector encoder, decoder and box heads
are bypassed.

The PHM scalar-image head receives extraction latency alongside pooled image features.
The dependence of latency on the computing environment limits attribution of the
regressor's performance to image content. The target measures candidate area in
photographs. Earlier target versions selected the maximum across views, while the
retained version uses a fixed view for each tooth. The top-three run summary aggregates
the tooth values. Independent physical defect masks were unavailable to validate
localisation.

\subsection{Shared RGB Preprocessing}
\label{sec:rtdetr-preprocessing}
For the \texttt{PekingU/rtdetr\_r18vd} checkpoint at revision
\nolinkurl{ac77a11ff0170a41b771c03264987f8ce2b0d753}, the retained processor uses
bilinear resizing to $640\times640$ and RGB rescaling by $1/255$.
ImageNet mean/standard-deviation normalisation is disabled
(\texttt{do\_normalize=false}), even though the configuration lists mean and standard-deviation
values. This preprocessing applies to the specified REASSEMBLE and ViFailback
paths; the PHM Ultralytics detector uses its separate implementation.\ev{Pinned processor
and checkpoint identity}{\EvidenceRef{E146}{}, \EvidenceRef{E149}{}}\ev{ViFailback
processor reuse}{\EvidenceRef{E092}{}, \EvidenceRef{E095}{}}

\section{Temporal Representations}
The PHM implementation defines a compact channel-independent PatchTST-style regressor
with default patch length sixteen, stride eight, hidden width 32, four attention heads,
two encoder layers, feed-forward width 64, dropout 0.1 and a 64-unit head. The
configured input width remains run-specific. A scalar output and the optional
sorted-profile head answer different target questions. The latter uses a base and
positive softplus increments to enforce ascending order; this guarantees an
order-statistic profile but discards tooth identity. The monotonicity does not
demonstrate that the estimated profile matches actual damage.\ev{PHM temporal
configuration}{\EvidenceRef{E029}{}}

The REASSEMBLE temporal branch follows a separate configuration: patch length 32, stride
sixteen, width 32, two layers, four heads, feed-forward width 64 and dropout 0.1. It
begins without pretrained temporal weights. Mean patch pooling and channel concatenation
feed a 128-dimensional projection and binary head, yielding 109,505 parameters in the
configured implementation. With internal scaling and input masking disabled, the observed-value mask does not
suppress missing tokens in attention. The encoder processes 31 patches per channel;
residual missing values are filled using training-only normalisation before encoding.
These are properties of the configured implementation, not guarantees implied by the
model family \parencite{Nie2023PatchTST}.\ev{REASSEMBLE temporal
recipe}{\EvidenceRef{E051}{}}

\section{Heads, Fitting and Calibration}
The original visual head is a new 896--128--1 network with ReLU and dropout 0.1,
containing 114,945 trainable parameters above a frozen backbone. Original visual epoch
candidates are five and ten; temporal candidates are three and six. Both use AdamW with
learning rate 0.001, weight decay 0.001 and batch size 128. These candidates constitute
the recorded selection procedure, not a broad optimiser search. Logistic regression
separately compares $C\in\{0.1,1,10\}$ under inner-fold selection.\ev{Original branch
configuration}{\EvidenceRef{E051}{}}

Branch calibration used an unweighted sigmoid logistic fit to the selected inner
out-of-fold scores. The outer assessment labels were withheld from calibration. The
original procedure selected a threshold between 0.01 and 0.99 by maximising inner
out-of-fold balanced accuracy, resolving ties toward 0.5. Later studies used a fixed
threshold of 0.5. Comparisons between those procedures therefore combine the effects of
model and threshold changes.

\section{Fusion Mechanisms and Attribution}
In the original REASSEMBLE procedure, uniform fusion averages the two calibrated branch
probabilities. F2 fits an unweighted L2 logistic model to those probabilities, with
$C=1$ and a new intercept and two coefficients. It is neither an F-beta metric nor a
convex weighted average. Its
coefficients depend on score scaling and are not causal percentages of modality
importance.

Feature fusion F3 combines a 128-dimensional visual projection with the 220 statistical
numerical features through separate projections, normalisation and a fitted classifier.
Its recorded candidates use hidden widths 32 or 64 and ten or twenty epochs. Context
gates instead allocate weights over branch logits. Their candidates have hidden width
eight or sixteen and forty epochs. Action-only, quality-only and combined contexts
remain different models, as do a static gate with context held constant and a gate
exposed to corrupted OOF records.\ev{Fusion and gate
configurations}{\EvidenceRef{E052}{}}

With one branch available, the explicitly masked convex-logit gate gives that branch a
weight of one. The resulting prediction follows from the specified masking rule. When
both branches are absent, some implementations reject the input and the external
evaluation records abstention and coverage.

The later anchored adjustment is another distinct mechanism. It modifies frozen F2
probability coefficients by bounded context-dependent factors, with a zero-initialised
adaptive output. It is not the same convex logit gate and need not have the same
calibration or threshold policy. Appropriate static and exposure-matched controls are
necessary to identify what, if anything, adaptation adds.\ev{Anchored
implementation}{\EvidenceRef{E061}{}}

\section{External Adaptations}
MSM and HANDOVER use the shared frozen visual backbone with a newly trained head and a
dataset-specific numerical classifier.\ev{External model register}{\EvidenceRef{E106}{}}

MSM fixes the HGB sensor control at 100 iterations, learning rate 0.1, fifteen maximum
leaf nodes, minimum leaf size ten and L2 regularisation 1.0, with early stopping
disabled. Training-only inverse-class-frequency sample weights balance its loss; these
settings were not selected by a hyperparameter search. The logistic sensor branch uses
training-only standardisation, balanced L2 logistic regression with $C=1$, the L-BFGS
solver and at most 2,000 iterations. Original F2 uses the same solver, iteration limit
and $C$, but is unweighted and fits a new intercept and two coefficients to clean
group-OOF branch probabilities. The primary evaluation fits no separate calibrator and
uses a fixed threshold of 0.5.\ev{Fixed MSM model and fusion configuration}{\EvidenceRef{E066}{},
\EvidenceRef{E147}{}, \EvidenceRef{E148}{}}

ViFailback adds pooling across three camera views. Its matched study compares frozen
generic and transferred CNN weights with a fresh target head in each case. Scalers,
sigmoid calibrators and out-of-fold fusion coefficients were fitted to target data
where specified. FROZEN20 and DONOR20 provide complete three-seed, three-outer-fold
assessments. FULL20 covers seed~42, outer fold~0 and its three inner fits; FULL100 is a
single-role pilot alongside FROZEN100 and DONOR100. Missing, degraded and temporally
shifted input conditions remain unmeasured. Table~\ref{tab:vif-status} gives the exact
fit counts and assessment scope.\ev{Completed ViFailback comparison scope}{\EvidenceRef{E113}{},
\EvidenceRef{E114}{}}
